\documentclass[10pt]{extarticle}
\usepackage[a4paper,margin=1.5cm]{geometry}
\usepackage[T1]{fontenc}
\usepackage{amsmath,amssymb}
\usepackage{graphicx}
\usepackage[export]{adjustbox}
\usepackage{multicol}
\usepackage{enumitem}
\usepackage{booktabs,array,tabularx}
\setlength{\columnsep}{0.7cm}
\setlength{\parindent}{0pt}
\setlength{\parskip}{3pt}
\pagestyle{empty}
\makeatletter
\newcommand{\qfigdef}[4]{\@namedef{qf@file@#1}{#2}\@namedef{qf@w@#1}{#3}\@namedef{qf@wide@#1}{#4}}
\newcommand{\qf@check}[1]{\@ifundefined{qf@file@#1}{\PackageError{score2}{Unknown figure id #1}{}}{}}
\newcommand{\qf@img}[4][]{\includegraphics[width=#2\dimexpr\@nameuse{qf@w@#4}pt\relax,max width=#3,max height=0.45\textheight,#1]{\@nameuse{qf@file@#4}}}
\newif\ifqf@full
\newif\ifqf@boxed
\newsavebox{\qf@box}
\newcommand{\qf@open}{\global\qf@fulltrue\par\global\qf@boxedtrue
  \setbox\qf@box\vbox\bgroup\hsize\textwidth\linewidth\textwidth}
\newcommand{\qf@emit}[1]{\par\egroup\global\qf@boxedfalse
  \ifdim\dimexpr\ht\qf@box+\dp\qf@box\relax<#1\textheight\noindent\usebox{\qf@box}\par\else\unvbox\qf@box\fi}
\newcommand{\qfullsplit}{\ifqf@boxed\qf@emit{0.85}\fi}
\newcommand{\qf@close}{\ifqf@boxed\qf@emit{0.45}\fi\global\qf@fullfalse}
\newcommand{\qfullbegin}{\end{multicols}\qf@open}
\newcommand{\qfullend}{\qf@close\begin{multicols}{2}}
\newcommand{\qfullstart}{\qf@open}
\newcommand{\qfulldone}{\qf@close}
\newcommand{\qfig}[2][1]{\qf@check{#2}%
  \ifqf@full\qf@blk{#1}{#2}\else
  \ifnum\@nameuse{qf@wide@#2}=1 \expandafter\@firstoftwo\else\expandafter\@secondoftwo\fi
  {\end{multicols}\begin{center}\qf@img{#1}{\textwidth}{#2}\end{center}\begin{multicols}{2}}%
  {\qf@blk{#1}{#2}}\fi}
\newcommand{\qf@blk}[2]{\par\smallskip{\centering\qf@img{#1}{\linewidth}{#2}\par}\smallskip}
\newcommand{\qfigi}[2][1]{\qf@check{#2}\qf@img[valign=c]{#1}{\linewidth}{#2}}
\makeatother
\newcommand{\papersec}[1]{\par\bigskip{\centering\Large\bfseries PAPER-#1\par}\medskip}
\newcommand{\partsec}[2]{\par\medskip\noindent{\large\bfseries PART-#1: #2}\par\smallskip}
\newcommand{\sect}[1]{\par\smallskip\noindent\textbf{#1}\par}
\newcommand{\qn}[1]{\par\smallskip\noindent\textbf{#1.}~}
\newcommand{\paraq}[2]{\par\smallskip\noindent\textbf{Paragraph for Questions #1 and #2}\par}
\newenvironment{qoptions}{\begin{itemize}[leftmargin=1.8em,itemsep=1pt,topsep=2pt,label={}]}{\end{itemize}}
\newcommand{\opt}[2]{\item[(#1)] #2}
% figure registry (generated by score2_tools.py finalize)
\qfigdef{p1_phy_s1_q01_1}{Score-II - Test 14 (06 Feb 2024)-img/p1_phy_s1_q01_1.png}{156.8}{0}
\qfigdef{p1_phy_s1_q02_1}{Score-II - Test 14 (06 Feb 2024)-img/p1_phy_s1_q02_1.png}{56.5}{0}
\qfigdef{p1_phy_s1_q03_1}{Score-II - Test 14 (06 Feb 2024)-img/p1_phy_s1_q03_1.png}{124.9}{0}
\qfigdef{p1_phy_s1_q04_1}{Score-II - Test 14 (06 Feb 2024)-img/p1_phy_s1_q04_1.png}{140.6}{0}
\qfigdef{p1_phy_s2_q05_1}{Score-II - Test 14 (06 Feb 2024)-img/p1_phy_s2_q05_1.png}{103.0}{0}
\qfigdef{p1_phy_s2_q06_1}{Score-II - Test 14 (06 Feb 2024)-img/p1_phy_s2_q06_1.png}{89.7}{0}
\qfigdef{p1_phy_s2_q07_1}{Score-II - Test 14 (06 Feb 2024)-img/p1_phy_s2_q07_1.png}{129.6}{0}
\qfigdef{p1_phy_s2_q08_1}{Score-II - Test 14 (06 Feb 2024)-img/p1_phy_s2_q08_1.png}{148.2}{0}
\qfigdef{p1_phy_s2_q09_1}{Score-II - Test 14 (06 Feb 2024)-img/p1_phy_s2_q09_1.png}{76.5}{0}
\qfigdef{p1_phy_s2_q10_1}{Score-II - Test 14 (06 Feb 2024)-img/p1_phy_s2_q10_1.png}{154.5}{0}
\qfigdef{p1_phy_s3_para1_2_1}{Score-II - Test 14 (06 Feb 2024)-img/p1_phy_s3_para1_2_1.png}{119.3}{0}
\qfigdef{p1_phy_s3_para3_4_1}{Score-II - Test 14 (06 Feb 2024)-img/p1_phy_s3_para3_4_1.png}{232.4}{0}
\qfigdef{p1_phy_s4_q09_1}{Score-II - Test 14 (06 Feb 2024)-img/p1_phy_s4_q09_1.png}{82.0}{0}
\qfigdef{p1_chem_s1_q04_1}{Score-II - Test 14 (06 Feb 2024)-img/p1_chem_s1_q04_1.png}{76.5}{0}
\qfigdef{p1_chem_s1_q04_A}{Score-II - Test 14 (06 Feb 2024)-img/p1_chem_s1_q04_A.png}{71.1}{0}
\qfigdef{p1_chem_s1_q04_B}{Score-II - Test 14 (06 Feb 2024)-img/p1_chem_s1_q04_B.png}{82.0}{0}
\qfigdef{p1_chem_s1_q04_C}{Score-II - Test 14 (06 Feb 2024)-img/p1_chem_s1_q04_C.png}{59.2}{0}
\qfigdef{p1_chem_s1_q04_D}{Score-II - Test 14 (06 Feb 2024)-img/p1_chem_s1_q04_D.png}{74.0}{0}
\qfigdef{p1_chem_s2_q09_A}{Score-II - Test 14 (06 Feb 2024)-img/p1_chem_s2_q09_A.png}{79.1}{0}
\qfigdef{p1_chem_s2_q09_2}{Score-II - Test 14 (06 Feb 2024)-img/p1_chem_s2_q09_2.png}{69.4}{0}
\qfigdef{p1_chem_s2_q09_B}{Score-II - Test 14 (06 Feb 2024)-img/p1_chem_s2_q09_B.png}{84.5}{0}
\qfigdef{p1_chem_s2_q09_4}{Score-II - Test 14 (06 Feb 2024)-img/p1_chem_s2_q09_4.png}{60.4}{0}
\qfigdef{p1_chem_s2_q09_C}{Score-II - Test 14 (06 Feb 2024)-img/p1_chem_s2_q09_C.png}{135.6}{0}
\qfigdef{p1_chem_s2_q09_D}{Score-II - Test 14 (06 Feb 2024)-img/p1_chem_s2_q09_D.png}{73.0}{0}
\qfigdef{p1_chem_s2_q09_7}{Score-II - Test 14 (06 Feb 2024)-img/p1_chem_s2_q09_7.png}{65.2}{0}
\qfigdef{p1_chem_s2_q10_1}{Score-II - Test 14 (06 Feb 2024)-img/p1_chem_s2_q10_1.png}{269.4}{1}
\qfigdef{p1_chem_s3_q04_1}{Score-II - Test 14 (06 Feb 2024)-img/p1_chem_s3_q04_1.png}{316.2}{1}
\qfigdef{p1_chem_s4_q09_1}{Score-II - Test 14 (06 Feb 2024)-img/p1_chem_s4_q09_1.png}{61.8}{0}
\begin{document}
\begin{center}{\Large\bfseries JEE (Main+Advanced) Score-II --- Test 14 (06 Feb 2024)}\end{center}
\begin{multicols}{2}
\partsec{1}{PHYSICS}

\sect{Section-I (i): Single Correct}

\qn{1} Two rods AC and BC are hinged at the end `C' and the ends A and B are moving along the surface with velocities $v_A=8$ m/s and $v_B=4\sqrt3$ m/s respectively as shown. the velocity of hinged point `C' is
\qfig{p1_phy_s1_q01_1}
\begin{qoptions}
\opt{A}{$\sqrt{13}$ m/s}
\opt{B}{$2\sqrt{13}$ m/s}
\opt{C}{$3\sqrt{13}$ m/s}
\opt{D}{$4\sqrt{13}$ m/s}
\end{qoptions}

\qn{2} A uniform spherical shell of mass M = 4.5 kg and radius R = 8.5 cm can rotate about a vertical axis on frictionless bearings. A massless cord passes around the equator of the shell, over a pulley of rotational inertia $I=3.0\times10^{-3}\ kg\cdot m^2$ and radius r = 5.0 cm, and is attached to a small object of mass m = 0.60 kg. There is no friction on the pulley's axle; the cord does not slip on the pulley. What is the speed of the object when it has fallen 92 cm after being released from rest? (g = 10 $m/s^2$)
\qfig{p1_phy_s1_q02_1}
\begin{qoptions}
\opt{A}{1.52 m/s}
\opt{B}{1.35 m/s}
\opt{C}{1.85 m/s}
\opt{D}{2.25 m/s}
\end{qoptions}

\qn{3} AB and CD are two uniform wires of same resistance per unit length and lengths 120 cm and 100 cm respectively. The cell of emf 6V is ideal and the other cell of emf E has internal resistance $2\Omega$. A length of 20 cm of wire CD is balanced by 40 cm length of wire AB. Find emf E in volts, if the reading of ideal ammeter is 2A. All other connecting wires have negligible resistance.
\qfig{p1_phy_s1_q03_1}
\begin{qoptions}
\opt{A}{20 V}
\opt{B}{16 V}
\opt{C}{14 V}
\opt{D}{None of these}
\end{qoptions}

\qn{4} A thin wire of length 2a is placed along x-axis in the x-y plane. The wire carries a uniform charge density $\lambda$ as shown in figure. A point P is located at (2a, 0) and point S at (0, y). If the electric potential at points P and S are equal then
\qfig{p1_phy_s1_q04_1}
\begin{qoptions}
\opt{A}{$y=2a$}
\opt{B}{$y=a$}
\opt{C}{$y=\sqrt3a$}
\opt{D}{None of these}
\end{qoptions}

\sect{Section-I (ii): One or More Correct}

\qn{5} A horizontal syringe, as shown in Figure, contains an ideal fluid of density $\rho$. The cross-section area of plunger is A and the small orifice at the end of the cylinder has area a (A $\gg$ a). Select correct option(s).
\qfig{p1_phy_s2_q05_1}
\begin{qoptions}
\opt{A}{The velocity of outflow of liquid from the the syringe if the plunger moves with a constant velocity under the action of force F is $\sqrt{\dfrac{2F}{\rho A}}$.}
\opt{B}{If the volume of fluid is V then constant force F required to empty the cylinder in time t is $\dfrac{\rho AV^2}{2a^2t^2}$.}
\opt{C}{If the volume of fluid is V then work required to empty the cylinder by a constant force in time t is $\dfrac{\rho V^3}{2a^2t^2}$.}
\opt{D}{If the volume of fluid is V then constant force F required to empty the cylinder in time t is proportional to $t^2$.}
\end{qoptions}

\qn{6} A dipole consists of two equal and opposite charges separated by a small distance. One such a dipole with charges +q and $-q$ each, and mass m each connected by a massless rigid rod of length $\ell$, is rotating with angular speed $\omega_0$ about $-$ve z-axis in XY plane with its centre of mass at origin with zero initial velocity. There exists uniform magnetic field B in the region along $-$ve z-direction. Select correct option(s).
\qfig{p1_phy_s2_q06_1}
\begin{qoptions}
\opt{A}{The maximum speed of the centre of mass of the dipole during the motion is $\dfrac{BqL}{m}$.}
\opt{B}{If the magnitude of $\omega_0$ is greater than 2qB/m, the dipole will not stop rotating.}
\opt{C}{If the angular velocity becomes zero, speed of center of mass is maximum at that time.}
\opt{D}{The maximum speed of the centre of mass of the dipole during the motion is $\dfrac{\omega_0L}{2}$.}
\end{qoptions}

\qn{7} A charged particle (+q) is moving simple harmonically on the x-axis with its mean position at origin. Amplitude of the particle is A and its angular frequency is $\omega$. Then choose the correct option(s) :
\qfig{p1_phy_s2_q07_1}
\begin{qoptions}
\opt{A}{The magnitude of magnetic field at (2A, 0) will change periodically with period $2\pi/\omega$.}
\opt{B}{The maximum magnitude of the magnetic field at (0, A) is $\dfrac{\mu_0}{4\pi}\cdot\dfrac{q\omega}{A}$}
\opt{C}{The magnetic field at (A, A) at the moment the particle passes through $\left(\dfrac A2,0\right)$, will be $\dfrac{\sqrt3\cdot\mu_0q\omega}{5\sqrt5\pi A}$}
\opt{D}{The magnitudes of magnetic field at (0, A) and (0, $-$A) will be same at any time.}
\end{qoptions}

\qn{8} A source of sound S with natural frequency $f_0=1800$ Hz moves uniformly along a straight line AB separated from a stationary observer O by a distance $\ell=250$ m. The velocity of the source is equal to $\eta=0.80$ fraction of the velocity of sound. Answer the following questions.
\qfig{p1_phy_s2_q08_1}
\begin{qoptions}
\opt{A}{The frequency of sound received by the observer at the moment when the source gets closest to him is 5000 Hz}
\opt{B}{The frequency of sound received by the observer at the moment when the source gets closest to him is 2000 Hz}
\opt{C}{The distance between the source and the observer at the moment when the observer receives a frequency $f=f_0$ is nearly 320m}
\opt{D}{The distance between the source and the observer at the moment when the observer receives a frequency $f=f_0$ is nearly 480m}
\end{qoptions}

\qn{9} Figure shows part of two waves, going in opposite direction, each of same wavelength 40 cm and same amplitude and speed. Choose the correct options :
\qfig{p1_phy_s2_q09_1}
\begin{qoptions}
\opt{A}{The position of nodes are $x=15,35,55$ cm etc.}
\opt{B}{The position of nodes are $x=5,25,45$ cm etc.}
\opt{C}{The position of antinodes are $x=0,20,40$ cm etc.}
\opt{D}{The position of antinodes are $x=5,25,45$ cm etc.}
\end{qoptions}

\qn{10} In a Young's experiment the light source is at distance $\ell_1=20\ \mu m$ and $\ell_2=40\ \mu m$ from the slits. the light of wavelength $\lambda=500$ nm is incident on slits separated at a distance 10 $\mu m$. A screen is placed at a distance D = 2 m away from the slits as shown in figure.
\qfig{p1_phy_s2_q10_1}
\begin{qoptions}
\opt{A}{Number of maxima on screen is 39}
\opt{B}{Number of maxima on screen is 60}
\opt{C}{Order of maxima at point C is 40.}
\opt{D}{If light ray only between S \& $S_1$ passes through a medium of refractive index $\mu$, then for a specific value of $\mu$, zero order maxima is obtained at C.}
\end{qoptions}

\sect{Section-II (i): Paragraph Type}

\paraq{1}{2}
The insulated box shown in the given figure has an insulated partition which can slide without friction along the length of the box. Initially each of the two chambers of the box has 1 mol of a monatomic ideal gas ($\gamma=5/3$) at a pressure $P_0$, volume $V_0$ and temperature $T_0$. The chamber on the left is slowly heated by an electric heater so that its gas, pushing the partition, expands until the final pressure in both the chambers becomes $243P_0/32$. Determine
\qfig{p1_phy_s3_para1_2_1}

\qn{1} Final temperature of the gas in left chamber is $xT_0$ where x is

\qn{2} Work done by the gas in right chamber is $-yRT_0$ where y is - (R is universal gas constant)

\paraq{3}{4}
A convex lens of focal length 15 cm is split into two halves and the two halves are placed at a separation of 120 cm as shown in figure. Between the two halves of convex lens a plane mirror is placed horizontally and at a distance of 4 mm below the principal axis of the lens halves. An object AB of length 2 mm is placed at a distance of 20 cm from one half lens as shown in figure. Let final image is formed after following successive phenomena :(i) Refraction at left part of lens ; (ii) Reflection at the mirror and (iii) Refraction at right part of lens.
\qfig{p1_phy_s3_para3_4_1}

\qn{3} The distance of final image of point A form the principal axis is K mm, then find 12K.

\qn{4} If left lens is shifted downward by 1 mm then the distance of final image of point B from the x axis is K mm. Find the value of 6K.

\paraq{5}{6}
A hydrogen atom in ground state is moving with a kinetic energy of 30 eV. It collides with a deuterium atom in ground state at rest. The hydrogen atom is scattered at right angle to its original line of motion. Assume that energy of nth state in both the atoms is given by $E_n=-\dfrac{13.6}{n^2}$eV and the mass of deuterium is twice that of hydrogen.

\qn{5} If after collision one of the atoms is excited to n = 2 then what will be kinetic energy of deuterium aftern collision (give answer in eV)

\qn{6} What will be maximum possible kinetic energy of deuterium (given answer in eV)

\sect{Section-II (ii): Numerical Value}

\qn{7} Consider a rocket that is in deep space and at rest relative to an inertial reference frame. the rocket's engine is to be fired for a certain interval. What must be the rocket's mass ratio (ratio of initial to final mass) to the nearest integer over that interval if the rocket's original speed relative to the inertial frame is to be equal to 2.0 times the exhaust speed (speed of the exhaust products relative to the rocket)? (e = 2.7)

\qn{8} Light of wavelength $\lambda$ from a small 0.5 mW He-Ne laser source, used in school laboratory, shines from a spacecraft of mass 1000 kg. The time needed for spacecraft to reach a velocity of 1 km/sec. from rest is $P\times10^{17}$ second. Find value P.

\qn{9} A ring of radius r made of wire of density $\rho$ is rotated about a stationary vertical axis passing through its centre and perpendicular to the plane of the ring as shown in the figure. Determine the angular velocity (in rad/s) of ring at which the ring breaks. The wire breaks at tensile stress $\sigma$. Ignore gravity. take $\sigma/\rho=4$ (in SI units) and r = 1 m.
\qfig{p1_phy_s4_q09_1}

\partsec{2}{CHEMISTRY}

\sect{Section-I (i): Single Correct}

\qn{1} Sodium metal is produced industrially by the electrolysis of molten NaCl, while steel is produced by CO reduction of ores containing Fe(III) in a blast furnace. Tantalum is another useful metal, often used surgically to repair bone because of its high resistance to corrosion. Which of the following methods could be used industrially to produce tantalum from naturally occurring Ta(V) ?\\
$Na^+(aq)+e^-\rightleftharpoons Na(s)$\hfill $E^\circ=-2.73$ V\\
$Fe^{3+}(aq)+3e^-\rightleftharpoons Fe(s)$\hfill $E^\circ=+0.82$ V\\
$Ta_2O_5(s)+10H^+(aq)+10e^-\rightleftharpoons2Ta(s)+5H_2O(\ell)$\hfill $E^\circ=-0.75$ V\\
$CO_2(g)+2H^+(aq)+2e^-\rightleftharpoons CO(g)+H_2O(\ell)$\hfill $E^\circ=-0.10$ V
\begin{qoptions}
\opt{A}{Either reduction with sodium metal or reduction in a blast furnace.}
\opt{B}{Either electrolysis or reduction with sodium metal.}
\opt{C}{Electrolysis only.}
\opt{D}{Reduction in a blast furnace only.}
\end{qoptions}

\qn{2} In the given following, the \textbf{incorrect} statement is :
\begin{qoptions}
\opt{A}{$N_2O_3$ is pale blue liquid and is acidic oxide in nature}
\opt{B}{$NO_2$ (gas) is paramagnetic while in its solid state is a diamagnetic}
\opt{C}{Aqua regia dissolves, Pt and Au metal and liberate $O_2$ gas.}
\opt{D}{With excess of Chlorine gas, $NH_3$ forms explosive compound $NCl_3$}
\end{qoptions}

\qn{3} Read the following statements :\\
(I) In the electrolytic reduction of alumina dissolved in fused cryolite ($Na_3AlF_6$) mixed with fluorspar ($CaF_2$), cryolite and fluorspar are used to decrease the M.P. and increase the conductivity of the electrolyte\\
(II) Leaching is a method to dissolve $Al_2O_3$ from red bauxite.\\
(III) The conversion of PbS into Pb requires roasting and self reduction.\\
(IV) Sodium can be extracted by the electrolytic-reduction of aqueous sodium chloride.\\
Pick the correct set of statements.
\begin{qoptions}
\opt{A}{I, II only}
\opt{B}{I, III, IV only}
\opt{C}{II, III, IV only}
\opt{D}{I, II, III only}
\end{qoptions}

\qn{4} Compound A, $C_{11}H_{12}O$, which gave no Tollen's test, was treated with $LiAlH_4$ followed by dilute acid, to give compound B, which could be resolved into enantiomers. When optically active B was treated with $CrO_3$ in pyridine, an optically inactive sample of A was obtained. Heating A with hydrazine in base gave hydrocarbon C, which when heated with alkaline $KMnO_4$, gave carboxylic acid D.
\qfig{p1_chem_s1_q04_1}
What is the structure of compound `A':
\begin{qoptions}
\opt{A}{\qfigi{p1_chem_s1_q04_A}}
\opt{B}{\qfigi{p1_chem_s1_q04_B}}
\opt{C}{\qfigi{p1_chem_s1_q04_C}}
\opt{D}{\qfigi{p1_chem_s1_q04_D}}
\end{qoptions}

\sect{Section-I (ii): One or More Correct}

\qn{5} For the reaction\\
$Br^-(aq)+BrO_3^-(aq)+H^+(aq)\longrightarrow Br_2(aq)+H_2O(\ell)$\\
the rate expression was found to be $-d[BrO_3^-]/dt=k[Br^-][H^+]^2[BrO_3^-]$\\
Which of the following statement(s) is/are correct ?
\begin{qoptions}
\opt{A}{Molecularity of reaction is 4.}
\opt{B}{Dimensions of pseudo rate constant of the reaction in a buffer solution is $M^{-1}min^{-1}$.}
\opt{C}{Doubling the concentration of all the reactants at the same time will increase the reaction rate by a factor of 16.}
\opt{D}{Rate of decrease of $BrO_3^-$ and rate of formation of $Br_2$ are the same.}
\end{qoptions}

\qn{6} The solubility of any solid in water may be classified as\\
Molecular solubility : $AB(s)\rightleftharpoons^{aq.}AB(aq)$\\
Ionic solubility : $AB(s)\rightleftharpoons^{aq.}A^+(aq)+B^-(aq)$\\
For a substance AB(s) the following thermodynamics informations are available at 300 K in standard conditions :\\
Molecular solubility : $\Delta H^\circ=+20$ kcal/mol; $\Delta S^\circ=+40$ cal/K-mole\\
Ionic solubility : $\Delta H^\circ=-25$ kcal/mol; $\Delta S^\circ=-50$ kcal/K-mole.\\
Which of the following statement(s) are incorrect for the solubility of AB(s) in water at 300 K in standard conditions :
\begin{qoptions}
\opt{A}{The salt will form a molecular solution}
\opt{B}{The salt will form an ionic solution}
\opt{C}{The salt will form both ionic as well as molecular solutions}
\opt{D}{The salt is completely insoluble in water.}
\end{qoptions}

\qn{7} Choose the \textbf{incorrect} statement from the following :
\begin{qoptions}
\opt{A}{H.O.M.O for $N_2$ molecule is $\pi_b$ M.O.}
\opt{B}{L.U.M.O for $N_2$ ion is $\pi^*$ M.O.}
\opt{C}{H.O.M.O for CO molecule is N.B.M.O from C-atom}
\opt{D}{Among halogen pale yellow colour of $F_2$ is due to $\sigma^*$ (H.O.M.O) to $\pi^*$ (L.U.M.O) transition}
\end{qoptions}

\qn{8} Which of the following statement(s) is/are \textbf{CORRECT} for complex (A).\\
$4FeCl_3+3K_4[Fe(CN)_6]\longrightarrow(A)+12KCl$
\begin{qoptions}
\opt{A}{Complex (A) is iron(II) hexacyanidoferrate(III)}
\opt{B}{E.A.N. of central metal atom of (A) complex is 36}
\opt{C}{Average oxidation number of Fe = $+2.50$}
\opt{D}{Total number of counter ions \& complex ions are = 7}
\end{qoptions}

\qn{9} Select option(s) with correct major product?\\
\begin{qoptions}
\opt{A}{\qfigi{p1_chem_s2_q09_A} $\xrightarrow[\text{(ii)}H_3O^+]{\text{(i)}Br_2/KOH}$ \qfigi{p1_chem_s2_q09_2}}
\opt{B}{\qfigi{p1_chem_s2_q09_B} $\xrightarrow{H_2SO_4}$ \qfigi{p1_chem_s2_q09_4}}
\opt{C}{\qfigi{p1_chem_s2_q09_C}}
\opt{D}{\qfigi{p1_chem_s2_q09_D} $\xrightarrow[\text{2. }\Delta]{\text{1. Base}}$ (after $O_3/Me_2S$) \qfigi{p1_chem_s2_q09_7}}
\end{qoptions}

\qfullbegin
\qn{10} \qfig{p1_chem_s2_q10_1}
\qfullsplit
Which of the following statement is/are correct :
\begin{qoptions}
\opt{A}{S is a non-narcotic analgesic which is also used in preventing platelet coagulation.}
\opt{B}{By the reaction of Y with formaldehyde in the presence of acid or base catalyst a elastomer polymer is obtained.}
\opt{C}{During the formation of P from Y, formation of a new C--C bond will takes place.}
\opt{D}{R undergoes Cannizzaro reaction and Z undergoes aldol condensation.}
\end{qoptions}
\qfullend

\sect{Section-II (i): Paragraph Type}

\paraq{1}{2}
Depending on the pH of the solution, $Mg^{2+}$ ions may hydrolyze to water soluble $MgOH^+$ species or precipitate as $Mg(OH)_2$.\\
$Mg(OH)_2(s)\rightleftharpoons Mg^{2+}(aq)+2OH^-(aq)$\hfill $K_{sp}=6\times10^{-12}$\\
$Mg^{2+}(aq)+OH^-(aq)\rightleftharpoons MgOH^+(aq)$\hfill $K_1=4\times10^2$\\
Consider a 10 litre solution of $2\times10^{-4}$ M $MgCl_2$ prepared in a buffer of pH = 11.\\
\relax[Atomic mass : Mg = 24, Cl = 35.5, O = 16, H = 1]

\qn{1} Determine \% $Mg^{2+}$ by mole (relative to all soluble Mg species) in this solution.

\qn{2} Calculate the mass of $Mg(OH)_2$ (in milli-grams) that will precipitate out from this solution.

\paraq{3}{4}
An aqueous blue coloured solution of a transition metal sulphate reacts with $H_2S$ in acidic medium to give a black precipitate A, which is insoluble in warm aqueous solution of KOH. The blue solution on treatment with KI in weakly acidic medium, turns yellow / brown solution due to the formations of compound C and produces a white precipitate B.\\
(given Atomic mass of Cu = 63.5, S = 32, O = 16, N = 14, I = 127, K = 39, H = 1)

\qn{3} Number of atoms in per formula unit of C.

\qn{4} "A" when dissolve in dilute nitric acid white cloudy appearance produce in test tube, then how many grams of dilute $HNO_3$ is required to dissolve 1 mole of "A" ?

\qfullbegin
\paraq{5}{6}
\qfig{p1_chem_s3_q04_1}
\qfullend

\qn{5} How many stereoisomers are possible of product `H'?

\qn{6} Degree of unsaturation of product `O' is :

\sect{Section-II (ii): Numerical Value}

\qn{7} A mixture of gases containing mainly carbon monoxide and hydrogen is produced by the reaction of alkene oxygen and steam as :\\
$CH_4(g)+\dfrac12O_2(g)\longrightarrow CO(g)+2H_2(g)$; $\Delta H=-36$ kJ/mol\\
$CH_4(g)+H_2O(g)\longrightarrow CO(g)+3H_2(g)$; $\Delta H=216$ kJ/mol\\
What wt. (in gm) of $CH_4$ gas required to produce 60 gm $H_2$ gas if there is no enthalpy change during above reaction.

\qn{8} $Mn^{+2}+Pb_3O_4\xrightarrow{H^+}$\\
$Mn^{+2}+PbO_2\xrightarrow{H^+}$\\
$Mn^{+2}+S_2O_8^{-2}\longrightarrow$\\
$Mn^{+2}+IO_4^-\xrightarrow{H^+}$\\
$Mn^{+2}+BiO_3^-\xrightarrow{H^+}$\\
$Mn^{+2}+MnO_4^-\longrightarrow$\\
$Mn^{+2}+OH^-\xrightarrow{\text{atomospheric oxidation}}$\\
$Mn^{+2}+OH^-\xrightarrow{H_2O_2}$\\
Number of reactions in which $MnO_4^-$ can be formed.

\qn{9} General structure of amino acid \qfigi{p1_chem_s4_q09_1}
\begin{center}
\begin{tabular}{|c|l|c|}
\hline
 & Name & Isoelectric point ($P^I$) \\
\hline
1 & Glycine & 5.97 \\
\hline
2 & Isoleucine & 6.02 \\
\hline
3 & Phenylalanine & 5.48 \\
\hline
4 & Aspartic acid & 2.76 \\
\hline
5 & Asparagine & 5.41 \\
\hline
6 & Arginine & 10.76 \\
\hline
7 & Lysine & 9.82 \\
\hline
8 & Serine & 5.68 \\
\hline
\end{tabular}
\end{center}
Referring the above table and find the sum of $P^I$ (isoelectric point) values of amino acids that satisfy each of the following criteria :\\
(a) The amino acid whose conjugated acid has two identical resonating structures.\\
(b) The amino acid that can exist as diastereomer.\\
(c) The amino acid that has zero optical rotation under all conditions\\
(d) The amino acid that are converted into other amino acid on treatment with conc. hot aq. NaOH solution followed by neutrzlization.

\partsec{3}{MATHEMATICS}

\sect{Section-I (i): Single Correct}

\qn{1} Let f \& g are two real valued continuous functions and let $\displaystyle\int g(x)\,dx=f^{-1}(x)$ and $f(x)=x^3+x+\sin\pi x+2$ then the value of $\displaystyle\int_2^4xg(x)$ is
\begin{qoptions}
\opt{A}{$\dfrac{-11}{4}-\dfrac2\pi$}
\opt{B}{$\dfrac{11}{4}-\dfrac2\pi$}
\opt{C}{$\dfrac{-11}{4}+\dfrac2\pi$}
\opt{D}{$\dfrac{11}{4}+\dfrac2\pi$}
\end{qoptions}

\qn{2} In a bag there are 10 black \& 10 white balls. A ball is drawn at random \& 5 extra balls of same color as of drawn ball are added in the bag along with the drawn ball. Now another ball is drawn and replaced back in the bag but 4 balls of same color as of drawn ball are removed from the bag. Again a ball is drawn and found to be white, then the probability that the second drawn ball was black, is
\begin{qoptions}
\opt{A}{$\dfrac47$}
\opt{B}{$\dfrac37$}
\opt{C}{$\dfrac27$}
\opt{D}{$\dfrac17$}
\end{qoptions}

\qn{3} Let A be a square matrix of order 3 satisfies $A^3-6A^2+7A-8I=O$ \& $B=A-2I$\\
$|adj(I-2A^{-1})|=K$, then the value of K, is (where I is an identity matrix of order `3')
\begin{qoptions}
\opt{A}{$\dfrac{25}{4}$}
\opt{B}{$\dfrac{25}{16}$}
\opt{C}{$\dfrac{25}{8}$}
\opt{D}{$\dfrac{25}{32}$}
\end{qoptions}

\qn{4} If $x^2+y^2+z^2=1$, x, y, z $\in R$ and maximum value of $(2x-y)^2+(3y-2z)^2+(z-3x)^2$ is $\lambda$, then $4\lambda$ is
\begin{qoptions}
\opt{A}{16}
\opt{B}{28}
\opt{C}{48}
\opt{D}{56}
\end{qoptions}

\sect{Section-I (ii): One or More Correct}

\qn{5} If $f(k)=\displaystyle\sum_{r=1}^k\dfrac1r$ and $\displaystyle\sum_{r=1}^{2017}f(r)=a\,f(b)+c$ (where a, b, c $\in$ I) then which of the following can be true ?
\begin{qoptions}
\opt{A}{b = c}
\opt{B}{a $-$ b = 1}
\opt{C}{b + c = 0}
\opt{D}{a + c = 0}
\end{qoptions}

\qn{6} A parabola S = 0 has its vertex at $(-9,3)$ and it touches the x-axis at the origin then equation of axis of symmetry of the parabola can be
\begin{qoptions}
\opt{A}{$x-y+12=0$}
\opt{B}{$x-2y+15=0$}
\opt{C}{$2x-y+21=0$}
\opt{D}{$x+y+6=0$}
\end{qoptions}

\qn{7} If $n_1$ = number of outcomes if 2 alike and 4 different dice are thrown, $n_2$ = number of outcomes if 6 alike dice are thrown and $n_3$ = number of divisors of $n_1$ of the form $4k+1$, $k\in N$, then
\begin{qoptions}
\opt{A}{$n_1+n_2=27678$}
\opt{B}{$n_1+n_2=27578$}
\opt{C}{$n_3=5$}
\opt{D}{$n_3=6$}
\end{qoptions}

\qn{8} Let A be the area bounded by the curve $f(x)=\dfrac{2}{1+x^4}$, x-axis, y-axis and ordinate x = 1, then-
\begin{qoptions}
\opt{A}{$A>\dfrac\pi2$}
\opt{B}{$A>2\ell n2$}
\opt{C}{$0<A<\dfrac\pi2$}
\opt{D}{$\dfrac\pi2<A<2\ell n2$}
\end{qoptions}

\qn{9} Let $S_1$ and $S_2$ be two circles touching each other externally at R. Let $\ell_1$ be a line which is tangent to $S_2$ at P and passing through centre $C_1$ of circle $S_1$. Similarly let $\ell_2$ be a line which is tangent to $S_1$ at Q and passing through centre $C_2$ of circle $S_2$. $\ell_1$ and $\ell_2$ are not parallel and intersect at K. If KP = KQ, then
\begin{qoptions}
\opt{A}{P and Q lie on opposite sides of the line joining $C_1$ and $C_2$}
\opt{B}{P and Q lie an same side of the line joining $C_1$ and $C_2$}
\opt{C}{Triangle PQR is an equilateral triangle}
\opt{D}{Triangle PQR is not necessarily an equilateral triangle}
\end{qoptions}

\qn{10} Let $\displaystyle\lim_{x\to\infty}\sqrt{x^4+Ax^3+Ax^2+x-1}-\sqrt{Cx^4+2x^3+7x^2-4x+1}=L$, then
\begin{qoptions}
\opt{A}{$C\ne1\Rightarrow$ L does not exist}
\opt{B}{$C=1,A\ne2\Rightarrow$ L does not exist}
\opt{C}{$C=1,A=2\Rightarrow L=-\dfrac52$}
\opt{D}{$A=C\Rightarrow$ L does not exist}
\end{qoptions}

\sect{Section-II (i): Paragraph Type}

\paraq{1}{2}
Let f be a differentiable function satisfying\\
$\sqrt[3]{f(x+y)}=\sqrt[3]{f(x)}+\sqrt[3]{f(y)}+1\ \forall\ x,y\in R$ and $f'(0)=3$

\qn{1} If $h(x)=f(x)-x^3$, then number of point(s) where $y=h(|x|)$ is non-derivable is (are):

\qn{2} If $x_0$ is solution of the equation $f(x)=f^{-1}(x)$, then $\left[\cos^{-1}(\cos2x_0)+4\tan^{-1}\left(\tan\dfrac{x_0}{2}\right)\right]$ is equal to (where $[*]$ denotes greatest integer function)

\paraq{3}{4}
Let A be a variable point on locus of feet of perpendicular drawn from focus upon any tangent to the curve $|z-2|+|z+2|=6$ and B be a variable point on $(1-i)z+(1+i)\bar z=10\sqrt2$, then

\qn{3} Minimum value of AB is

\qn{4} Let locus of variable point A is a curve $C_1$ and locus of variable point B is a curve $C_2$. If a variable circle touches both the curves $C_1$ and $C_2$ externally, then length of latus rectum of locus of centre of variable circle is :

\paraq{5}{6}
Given a triangle ABC (with usual nations) such that angles A, B, C $\in\left(0,\dfrac\pi2\right)$ and x, y, z are positive real numbers, then

\qn{5} If the maximum value of $(\cos A)^A(\cos B)^B(\cos C)^C$ is $\lambda^\pi$ then the value of `$8\lambda$' is -

\qn{6} If the minimum value of $\dfrac{\sqrt{x+x^3}}{x\sin P\cos^2P}+\dfrac{\sqrt{y+y^3}}{y\sin Q\cos^2Q}+\dfrac{\sqrt{z+z^3}}{z\sin R\cos^2R}$ (where P, Q, R $\in\left(0,\dfrac\pi2\right)$) is $\mu$ then $2(\mu)^2$ is

\sect{Section-II (ii): Numerical Value}

\qn{7} If the curve $y=y(x)$ which satisfy the differential equation $e^{\frac xy(1-y^2)}\left[y\dfrac{dx}{dy}-x\right]+\left[xy^2+y^3\dfrac{dx}{dy}\right]=0$ and $y(2)=1$ is given by $Ae^{\frac xy}+e^{xy}+Be^C=0$, then $(A+B+C)$ is

\qn{8} If $S_n=3+\dfrac{1+3+3^2}{3!}+\dfrac{1+3+3^2+3^3}{4!}\ldots$ upto n terms.\\
Then the value of $\left[\displaystyle\lim_{n\to\infty}S_n\right]$ is, (where [.] represent Greatest Integer Function)

\qn{9} Let $I_{(m,n)}=\displaystyle\int(\cos2x)^m(1-\tan^2x)^ndx$, such that\\
$4I_{(0,10)}+\dfrac{I_{(1,10)}}{5}=A\sin2x(1-\tan^2x)^{10}+BI_{(0,9)}+C$ (where C is constant of integration, m, n $\in$ W), then the value of $\dfrac{B}{8A}$ is

\end{multicols}
\end{document}
